\section{Diversidad de trace y estrategias de muestreo}

Para que trace mejorara el rendimiento sin dejar de ser interpretable, el equipo diseñó un pipeline de muestreo en varias etapas: primero el solver produce 64 candidate trace, y luego una puntuación heuristic (cost, length, constraint violation) decide cuáles trace vale más la pena conservar. Cada trace registra cost, exploration depth y plan segment; estos metadata también se usan para entrenar la supervisión de diversity.
\begin{enumerate}
    \item \textbf{Generación de trace:} se usa solver/Beam search para construir un trace global y registrar el intermediate cost.
    \item \textbf{Filtrado de trace:} se filtran los trace de baja calidad según violation, length y user goal.
    \item \textbf{Muestreo de trace:} se muestrean varios plan con el objetivo de la diversidad, para garantizar que se cubran distintas estrategias durante el entrenamiento y la inferencia.
    \item \textbf{Reordenamiento de trace:} la sequence que produce el Transformer se valida mediante un judge y luego se alinea con el original solver trace, para formar nuevos datos de entrenamiento.
\end{enumerate}
\begin{knowledgebox}{Los beneficios de la diversidad}
Esta estrategia de muestreo hace que el entrenamiento ya no dependa de un solo optimal plan, sino que el modelo aprenda la manera de pensar de «buscar varias soluciones factibles y elegir la mejor»; sobre todo en Soang, Maze y Nano optics, diversity mejora directamente el test-time performance.
\end{knowledgebox}

\subsection{Resumen de la sección}

Un muestreo diversificado de trace puede combinar la ventaja del deterministic solver con la flexibilidad del probabilistic planner, conservando la corrección del plan a la vez que aporta nuevos candidate, de manera que el modelo de search trace pueda seguir produciendo soluciones robustas en entornos unseen.

